\chapter{Related work}
\label{chap:related-work}

This thesis builds on four lines of work. The first applies \ac{RL} to the choice of channel and
transmit power for \ac{V2V} links. The second studies how several \ac{RL} agents can learn together,
usually through a server. The third removes the server and lets each node combine
what it learns only with its neighbors. The fourth asks which reward each agent of a team should learn
from. Figure~\ref{fig:rw-map} shows the four lines and what this thesis takes from each.

The chapter reviews the four lines in that order. For each one, it states the problem the line
addresses and what the key papers did.
Section~\ref{sec:rw-positioning} then compares the thesis with the closest works.

\begin{figure*}[tbp]
  \centering
  \begin{tikzpicture}[
      font=\linespread{1}\small,
      area/.style={draw=black!45, fill=black!4, rounded corners=2pt, align=left,
                   text width=7.3cm, inner sep=5pt},
      thesis/.style={draw=black!80, thick, fill=black!10, rounded corners=2pt, align=center,
                     text width=3.3cm, inner sep=6pt},
      feed/.style={-{Stealth[length=2mm]}, semithick, draw=black!60}]
    \node[area] (v2v) {\textbf{\ref{sec:rw-v2v}\enspace \acs{RL} for \acs{V2V} resource allocation}\\
      \emph{Takes:} the problem, the agent design, and the baseline scheme of
      \citet{li+22fedmarl}};
    \node[area, below=2.5mm of v2v] (frl) {\textbf{\ref{sec:rw-frl}\enspace Federated and multi-agent
      \acs{RL}}\\
      \emph{Takes:} agents that learn together by sharing their network weights};
    \node[area, below=2.5mm of frl] (dfl) {\textbf{\ref{sec:rw-dfl}\enspace Serverless decentralized
      learning}\\
      \emph{Takes:} gossip averaging among neighbors, with no server};
    \node[area, below=2.5mm of dfl] (rew) {\textbf{\ref{sec:rw-reward}\enspace Reward design in
      multi-agent systems}\\
      \emph{Takes:} a reward that judges each agent mainly by its own actions};
    \coordinate (mid) at ($(frl.east)!0.5!(dfl.east)$);
    \node[thesis, right=0.8cm of mid] (t) {\textbf{This thesis}\\[3pt]
      \acs{D3QN} agents choose channel and power for \acs{V2V} links, average their networks
      only with nearby vehicles, and learn from individual rewards};
    \draw[feed] (v2v.east) -- ($(t.west)+(0,0.45)$);
    \draw[feed] (frl.east) -- ($(t.west)+(0,0.15)$);
    \draw[feed] (dfl.east) -- ($(t.west)+(0,-0.15)$);
    \draw[feed] (rew.east) -- ($(t.west)+(0,-0.45)$);
  \end{tikzpicture}
  \caption[The four lines of related work]{The four lines of work reviewed in this chapter, numbered by
    section, and what this thesis takes from each.}
  \label{fig:rw-map}
\end{figure*}

\section{Reinforcement learning for V2V resource allocation}
% Trabajos que asignan canal y potencia a enlaces V2V con aprendizaje por refuerzo en lugar de optimización centralizada: Ye et al. (2019), Liang et al. (2019), Zhang et al. (2020) y Li et al. (2022).
\label{sec:rw-v2v}

\ac{V2V} links can reuse the channels that the cellular network has already assigned to its users, the
\acp{CUE} (Section~\ref{sec:bg-v2v}). Each \ac{V2V} pair must then choose a channel and a transmit
power. The choice must deliver its own safety messages on time without degrading the cellular links that
share the channel.

Classical solutions treat this as an optimization problem solved at a central point. They need the
instantaneous state of every channel, which is hard to track when vehicles move fast and costly to report
to the \ac{BS}. They are also myopic: they optimize the current decision, not long-term performance
\cite{li+22fedmarl}. For these reasons, a recent line of work learns the allocation with \ac{RL}
instead.

\citet{ye+19drl} proposed a decentralized scheme based on deep \ac{RL}. Each \ac{V2V} link is an agent
that picks its sub-band and transmit power from local information, without waiting for global channel
information. The goal is to meet the latency constraints of \ac{V2V} messages while limiting the
interference on the vehicle-to-infrastructure links. Decisions are local, but \citet{li+22fedmarl} point
out that schemes of this kind still train on data collected from all vehicles in one place, which costs
communication and raises privacy concerns.

\citet{liang+19spectrum} treated the same spectrum-sharing problem as a multi-agent one. Each \ac{V2V}
link is an agent with its own \ac{DQN}. All agents receive the same reward, which combines the sum
capacity of the vehicle-to-infrastructure links with the delivery of \ac{V2V} payloads. The authors chose
a common reward on purpose: it turns a problem in which links compete for spectrum into a fully
cooperative one. Learning is centralized and execution is distributed: while training, every agent has
access to this system-wide reward; once trained, each agent acts on its local observation only. This is a
form of \ac{CTDE} \cite{amato24cooperative}. To reduce the non-stationarity caused by agents that learn at
the same time, they added a ``fingerprint'' to each observation: the training iteration and the
exploration rate. Their agents improved both the capacity of the infrastructure links and the delivery
rate of \ac{V2V} payloads.

\citet{zhang+20v2xfed} brought federated learning into this problem. Each \ac{V2V} pair is a \ac{DQN}
agent that chooses its transmission mode and its radio resources. Its reward adds up the capacity of the
cellular users and the latency and reliability terms of all the pairs, so it is common to all of them. The
\ac{BS} groups the pairs into clusters with a graph algorithm whose edge weights are the large-scale gains
between pairs. The reason they give is that nearby pairs experience similar channels and observations.
Every few hundred subframes, the \ac{BS} averages the models of the pairs of each cluster and sends the
average back to them. The authors use federation mainly to obtain robust models that newly activated pairs
can use with little training; in their results, it performs about as well as their other learning
algorithms.

\citet{li+22fedmarl}, the work this thesis starts from, proposed a federated multi-agent deep \ac{RL}
scheme. Each \ac{V2V} pair is a \ac{D3QN} agent trained with \ac{PER} (Section~\ref{sec:pa-agent}). It
observes the gain of its own link on each channel, the gains from the \acp{CUE} to the \ac{BS}, how many
neighboring pairs used each channel in the previous slot, and the length of its queue. Neighbors are the
pairs within \SI{300}{\meter}, the maximum communication range of the \ac{3GPP} evaluation scenario
\cite{3gpp16v2x}. The agent then chooses a channel and a power level. Learning has two central
elements:
\begin{itemize}[nosep]
  \item \emph{A global reward.} The \ac{BS} computes a single reward at the end of each slot and
    broadcasts it to all agents.
  \item \emph{Model averaging at the \ac{BS}.} Periodically, the \ac{BS} collects the networks of all
    agents, averages them with \ac{FedAvg} \cite{mcmahan+17fedavg}, and sends the average back. The
    authors argue that this gives each agent implicit knowledge of the parts of the network it cannot
    observe.
\end{itemize}
They compared their scheme with a greedy rule, a cluster-based method, and a scheme in which a single
agent learns and the others copy its model. Their scheme achieved a higher \ac{CUE} sum rate, a higher
\ac{V2V} packet delivery rate, and a shorter average \ac{V2V} transmission delay than the greedy and
cluster-based methods, and it converged faster and to a higher reward than the single-agent scheme.

\section{Federated and multi-agent reinforcement learning}
% Revisiones y aplicaciones de FRL y MARL federado: Qi et al. (2021), Myakala y Kamatala (2023) y Cheruiyot et al. (2025).
\label{sec:rw-frl}

When several agents face related tasks, each one sees only part of the possible experience. The question
this line of work addresses is how the agents can learn from each other's experience. The idea is old:
\citet{tan93independent} found that agents that share their learned policies or their experience episodes
learn faster than independent agents, at the cost of communication. \Ac{FRL} recasts this sharing in the
terms of \ac{FL} (Section~\ref{sec:bg-fl}): agents keep their raw experience and exchange only model
parameters or gradients.

\citet{qi+21frl} reviewed \ac{FRL} and split it into two variants according to how the environment is
divided among agents:
\begin{itemize}[nosep]
  \item \emph{Horizontal \ac{FRL}.} Each agent interacts with its own, independent copy of the
    environment, and all agents have the same state and action spaces. They exchange parameters or
    gradients, either through a server or peer to peer, to learn faster from more experience.
  \item \emph{Vertical \ac{FRL}.} All agents act in one shared environment. Each one observes only part
    of the state, and the actions of all of them affect the dynamics and the rewards.
\end{itemize}
Among the application areas they review are communications and edge computing, and they use autonomous
driving as a typical example of horizontal \ac{FRL}.

The \ac{V2V} scheme of \citet{li+22fedmarl} does not fit neatly into either variant. Its agents have the
same observation and action spaces and exchange model parameters, as in horizontal \ac{FRL}. But they act
in one shared environment and interfere with each other, as in vertical \ac{FRL}. The models being
averaged therefore come from agents whose experiences are coupled, not independent.

Two recent surveys address the tension between federation and decentralization directly.
\citet{cheruiyot+25marlsurvey} distinguish three regimes of \ac{MARL}. In \ac{FRL}, a central server
relays the communication and averages the parameters of the agents. In decentralized \ac{MARL},
there is no server: agents exchange parameters peer to peer over a communication graph, which
the survey also calls a gossip network, and they may maximize common or individual rewards. In
noncooperative \ac{MARL}, each agent has its own reward and the agents play a game.
\citet{myakala+23scalable} review scalable, decentralized multi-agent \ac{FRL}. They name communication
efficiency, decentralized aggregation, and non-stationarity as the main challenges of scaling \ac{FRL} to
many agents, and review remedies such as selective communication, compression, and peer-to-peer
aggregation.

\section{Serverless decentralized learning}
% Consenso y gossip (Boyd et al., 2006; Olfati-Saber et al., 2007), gossip learning (Hegedus et al., 2019), D-PSGD (Lian et al., 2017), MARL con agentes en red (Zhang et al., 2018) y DFL con vecinos limitados (Ge et al., 2024).
\label{sec:rw-dfl}

\Ac{FL} avoids pooling raw data, but it keeps a server that collects and averages every model. All the
traffic of learning passes through that server. Decentralized learning removes it: each node exchanges
information only with its neighbors in a communication graph (Section~\ref{sec:bg-dfl}). The question is
whether such local exchanges can reach the same result as a central average, and how fast.

\paragraph{Consensus and gossip.} The theoretical basis is the study of consensus and gossip algorithms.
\citet{boyd+06gossip} analyzed randomized gossip: at each step, a randomly chosen node averages its value
with that of one neighbor. Repeating this drives every node to the network-wide average. They showed that
the time this takes depends on the structure of the graph, through the second largest eigenvalue of the
averaging matrix. One of the graphs they studied in detail is the geometric random graph, in which nodes
are connected when they are close to each other; it models wireless networks.
\citet{olfatisaber+07consensus} reviewed consensus algorithms for networked multi-agent systems, including
networks whose topology changes over time.

\paragraph{Training models without a server.} Three works carry this idea from averaging values to
training models:
\begin{itemize}
  \item \citet{lian+17dpsgd} proposed decentralized parallel stochastic gradient descent (D-PSGD). Each
    node takes gradient steps on its own data and averages its model with those of its neighbors. D-PSGD
    needs a total amount of computation comparable to the centralized version, while the busiest node
    communicates much less. On networks with low bandwidth or high latency, it ran up to an order of
    magnitude faster.
  \item \citet{hegedus+19gossip} compared gossip learning with \ac{FL}. In gossip learning, each node
    periodically sends its model to a random peer, which merges it with its own model and updates it on
    local data. Using a real trace of phones going online and offline, and the same compression for both
    approaches, they found gossip learning comparable to \ac{FL} overall, and better when the training
    data were spread evenly across nodes.
  \item \citet{ge+24dfl} proposed a \ac{DFL} scheme for heterogeneous Internet-of-Things devices. In each
    round, each device picks a fixed number of neighbors at random and receives their models. Instead of
    averaging parameters, it runs the neighbors' models on its own data, fuses their outputs, and
    distills the result into its own model. This lets devices with different model architectures learn
    from each other. In their experiments, the scheme needed fewer communication rounds than \ac{FedAvg}
    to reach a target accuracy. More neighbors improved accuracy but raised communication, and five
    neighbors gave the best balance.
\end{itemize}

\paragraph{Decentralized multi-agent \acs{RL}.} \citet{zhang+18networked} proposed a fully decentralized
actor-critic method for networked agents. Each agent receives a private reward that only it knows. Agents
exchange the parameters of their critics with their neighbors over a network that may change over time,
and combine them through a consensus step. Together, they maximize the average return of all agents.
Their convergence guarantees assume a linear approximation of the value function; they also report
simulations with nonlinear approximation.

\section{Reward design in multi-agent systems}
% Recompensa compartida frente a recompensas por agente: agentes independientes (Tan, 1993; Claus y Boutilier, 1998), recompensas de diferencia (Wolpert y Tumer, 2001), alineación y sensibilidad (Agogino y Tumer, 2008), VDN (Sunehag et al., 2018) y críticos locales y global en V2X (Parvini et al., 2023).
\label{sec:rw-reward}

\begin{table*}[tbp]
  \centering
  \linespread{1}\small
  \caption[Comparison with the closest works]{Comparison of this thesis with the closest works. The last
    column says whether the work measures the cost of communication. A dash (---) means that the aspect
    does not apply.}
  \label{tab:rw-positioning}
  \begin{tabularx}{\textwidth}{@{}>{\raggedright\arraybackslash}p{2.1cm}
      >{\raggedright\arraybackslash}X >{\raggedright\arraybackslash}X
      >{\raggedright\arraybackslash}X >{\raggedright\arraybackslash}p{1.6cm}@{}}
    \toprule
    Work & What agents share while learning & Who combines it & Reward each agent learns from
      & Comm.\ cost measured? \\
    \midrule
    \citet{liang+19spectrum} & Nothing; each agent trains its own \acs{DQN} & --- & Common to all
      & No \\
    \addlinespace
    \citet{zhang+20v2xfed} & \acs{DQN} weights & The \acs{BS}, within clusters of nearby pairs
      & Common to all & No \\
    \addlinespace
    \citet{li+22fedmarl} & \acs{D3QN} weights & The \acs{BS}, with \acs{FedAvg} & Common to all,
      computed by the \acs{BS} & No \\
    \addlinespace
    \citet{parvini+23aoi} & Observations and actions, sent to a shared critic & The roadside unit, in the
      shared critic & Individual, plus a team reward & No \\
    \addlinespace
    \citet{zhang+18networked} & Critic parameters & Each agent, with its neighbors (consensus)
      & Private to each agent & No \\
    \addlinespace
    \citet{lian+17dpsgd} & Weights of a supervised model & Each node, with its neighbors & --- & Yes \\
    \addlinespace
    \citet{hegedus+19gossip} & Weights of a supervised model & Each node, with one random peer at a
      time & --- & Yes \\
    \addlinespace
    \citet{ge+24dfl} & Supervised models, combined by distillation & Each device, with $S$ random
      neighbors & --- & Yes \\
    \midrule
    This thesis & \acs{D3QN} weights & Each vehicle, with its neighbors within \SI{300}{\meter} (all,
      or $S$ at random) & Individual, compared with the common one & Yes, as model transfers \\
    \bottomrule
  \end{tabularx}
\end{table*}


When agents cooperate, the simplest design gives all of them the same reward: a measure of how well the
whole system did. Its weakness is credit assignment. From a shared reward, an agent cannot tell how much of
the outcome came from its own action and how much from the actions of the others.
\citet{gronauer+22marlsurvey} list credit assignment among the central challenges of deep \ac{MARL}.

\paragraph{Independent and cooperative learners.} Early studies compared agents that learn on their own
with agents that cooperate. \citet{tan93independent} found that, on tasks that require joint action,
agents that work as partners end up outperforming independent ones, although they learn more slowly at
first. \citet{claus+98dynamics} compared, in cooperative games, Q-learners that ignore the other agents
with learners that model joint actions. They showed that both kinds may settle on a suboptimal
equilibrium, depending on the structure of the game and on how the agents explore.

\paragraph{Rewards that are aligned and learnable.} \citet{wolpert+01collectives} studied how to choose
the reward of each member of a collective so that members that maximize their own reward also make the
collective perform well. Giving every member the global reward, the so-called team game, is one option.
They derived alternatives that are easier for each member to learn and still aligned with the global goal.
One of them rewards each member with the difference between the global reward and the global reward
without that member's contribution, and it outperformed the team game in their experiments.
\citet{agogino+08rewards} made the trade-off explicit with two properties (Section~\ref{sec:bg-marl}):
\emph{alignment}, whether improving an agent's reward also improves the system's, and
\emph{learnability}, how sensitive the reward is to the agent's own actions compared with the actions of
the others. In a multi-rover domain, they showed that:
\begin{itemize}[nosep]
  \item the global reward is aligned but hard to learn, because each agent's effect is masked by the
    others;
  \item a purely local reward is easy to learn but may not be aligned;
  \item the difference reward has both properties.
\end{itemize}

\paragraph{Deep multi-agent \acs{RL}.} \citet{sunehag+18vdn} trained cooperative deep \ac{RL} agents from
a single team reward and observed two failures: spurious rewards, where an agent is rewarded for what a
teammate did, and the lazy agent, which stops learning because its exploration would lower the reward a
teammate already earns. They proposed value-decomposition networks (VDN): the value of the team is
written as a sum of per-agent values, and all of them are trained together, centrally, from the team
reward.

\paragraph{The reward in V2V resource allocation.} Both \citet{liang+19spectrum} and
\citet{li+22fedmarl} give all agents a common reward, and both justify it in the same way: the links
compete for channels, so with separate rewards they would act selfishly, and a common reward turns the
problem into a cooperative one. In the terms above, the price is learnability. The reward of
\citet{li+22fedmarl} adds one term per \ac{CUE} and two per \ac{V2V} pair, and each agent's action moves
only a few of them.

The remedies above do not fit a serverless setting easily. The difference reward needs the performance of
the system with and without each agent, VDN trains all agents together from the team reward, and the global
critic of \citet{parvini+23aoi} needs the observations and actions of every agent. All of them need a
central view of the system.

\section{Positioning of this thesis}
% Tabla que compara los trabajos más cercanos (qué se comparte, quién agrega, de qué recompensa aprende cada agente y si se mide la comunicación) y qué combina esta tesis que no aparece junto en ellos.
\label{sec:rw-positioning}

Table~\ref{tab:rw-positioning} compares the closest works along four questions: what the agents share
while they learn, who combines what is shared, what reward each agent learns from, and whether the cost of
communication is measured.

The table shows the gap. The \ac{V2V} works learn through a central element. All but one give every agent
the same reward, and the exception adds individual rewards to a centrally trained team critic. The
decentralized works combine models among neighbors, but they either train supervised models
or, in \ac{RL}, give guarantees only for linear critics on an abstract network and do not measure
communication. No work brings
together the three elements of this thesis: multi-agent \ac{RL} for \ac{V2V} resource allocation, model
averaging among neighbors without the \ac{BS}, and an individual reward per agent.

Against each of the closest works, the thesis differs as follows:
\begin{itemize}
  \item Compared with \citet{li+22fedmarl}, the thesis keeps the problem, the agents, and their
    observations and actions. It makes two changes: it removes the \ac{BS} from model averaging, and it
    replaces the common reward with an individual one. Within our implementation, everything else stays the
    same, so the differences between the methods can be attributed to these two changes. That
    implementation reproduces only part of the published behavior (Section~\ref{sec:res-replication}), so
    the comparison is with our implementation of the scheme, not with the published results. The gossip
    neighbors are the same \SI{300}{\meter} neighbors that the agents already use in their observation.
    Unlike \citet{li+22fedmarl}, the thesis also counts the model transfers that each form of aggregation
    needs and how they are spread across nodes.
  \item Compared with \citet{zhang+20v2xfed}, the thesis also averages models only among nearby pairs, and
    its results are consistent with their reason for doing so, that nearby pairs live similar situations,
    although they do not prove it (Section~\ref{sec:res-discussion}). Unlike them, there is no \ac{BS} in
    the averaging: each vehicle has its own group of neighbors, given by distance, and the vehicles
    average among themselves. Their agents also learn from a common reward.
  \item Compared with \citet{parvini+23aoi}, the thesis also gives each agent its own reward. Here,
    however, the individual reward replaces the common one instead of being added to it, no central critic
    is needed, and the two rewards are compared directly, in the same scenarios.
  \item Compared with \citet{zhang+18networked}, the thesis shares the combination of private rewards and
    consensus among neighbors. It differs in four ways: it uses deep Q-networks, outside the linear
    setting of their analysis; it works on a concrete communications problem; it derives the neighbor
    graph from the positions of the vehicles; and it measures the cost of communication empirically. In
    exchange, it offers no convergence guarantees.
  \item From \citet{ge+24dfl}, the thesis borrows the idea of capping at $S$ the number of neighbors each
    node combines with per round. Unlike them, it averages parameters instead of distilling outputs, since
    all agents share one architecture, and its nodes are \ac{RL} agents.
  \item With \citet{lian+17dpsgd} and \citet{hegedus+19gossip}, the thesis shares the idea of averaging
    models between peers. Here, however, the data come not from a fixed dataset but from agents that share
    one environment and interfere with each other.
\end{itemize}
